\labsection{2}{Supervised learning I: regression, classification, and the workflow}{sec:lab02-supervised-workflow}

\begin{labproject}{Two regression examples, two classification examples, and a saved model at the end}
\textbf{Goal.} Fit regression/classification pipelines, validate chronologically where needed, tune thresholds and hyperparameters, test once, and save models.

\medskip\noindent\textbf{Setup.}\par
\begin{lstlisting}[style=mlpython]
import json
from pathlib import Path

import joblib
import matplotlib.pyplot as plt
import numpy as np
import pandas as pd
import sklearn
from sklearn.compose import ColumnTransformer
from sklearn.ensemble import RandomForestClassifier, RandomForestRegressor
from sklearn.impute import SimpleImputer
from sklearn.linear_model import LinearRegression, LogisticRegression, Ridge
from sklearn.metrics import (
    ConfusionMatrixDisplay,
    PrecisionRecallDisplay,
    RocCurveDisplay,
    accuracy_score,
    average_precision_score,
    classification_report,
    confusion_matrix,
    mean_absolute_error,
    mean_squared_error,
    precision_score,
    r2_score,
    recall_score,
    roc_auc_score,
)
from sklearn.model_selection import GridSearchCV, StratifiedKFold, train_test_split
from sklearn.pipeline import Pipeline
from sklearn.preprocessing import OneHotEncoder, StandardScaler

PROJECT_ROOT = Path.cwd().parent
DATA = PROJECT_ROOT / "all_datasets"
MODELS = PROJECT_ROOT / "models"
MODELS.mkdir(exist_ok=True)
\end{lstlisting}

Use one helper for numerical imputation/scaling and categorical imputation/one-hot encoding.

\begin{lstlisting}[style=mlpython]
def make_preprocess(num_cols, cat_cols):
    return ColumnTransformer([
        ("num", Pipeline([
            ("impute", SimpleImputer(strategy="median")),
            ("scale", StandardScaler()),
        ]), num_cols),
        ("cat", Pipeline([
            ("impute", SimpleImputer(strategy="most_frequent")),
            ("onehot", OneHotEncoder(handle_unknown="ignore")),
        ]), cat_cols),
    ])
\end{lstlisting}

\medskip\noindent\textbf{Part A --- Regression: ordinary least squares versus ridge.}\par

Predict insurance \texttt{charges}. Compare OLS and ridge on the same features and split.

\begin{lstlisting}[style=mlpython]
insurance = pd.read_csv(DATA / "medical_cost_personal_dataset.csv")
print(insurance.head())

X = insurance.drop(columns="charges")
y = insurance["charges"]
num_cols = ["age", "bmi", "children"]
cat_cols = ["sex", "smoker", "region"]

X_train, X_valid, y_train, y_valid = train_test_split(
    X, y, test_size=0.20, random_state=42
)
\end{lstlisting}

MAE gives average dollar error; RMSE emphasizes extremes; $R^2$ compares squared error with the mean-label reference.

\begin{lstlisting}[style=mlpython]
fitted = {}
predictions = {}

for name, estimator in {
    "linear regression": LinearRegression(),
    "ridge": Ridge(alpha=10.0),
}.items():
    model = Pipeline([("preprocess", make_preprocess(num_cols, cat_cols)), ("model", estimator)])
    model.fit(X_train, y_train)
    fitted[name] = model
    predictions[name] = model.predict(X_valid)
    pred = predictions[name]
    print(f"\n{name}")
    print("MAE:", round(mean_absolute_error(y_valid, pred), 2))
    print("RMSE:", round(np.sqrt(mean_squared_error(y_valid, pred)), 2))
    print("R2:", round(r2_score(y_valid, pred), 3))
\end{lstlisting}

Coefficients give conditional prediction changes. Standardized numerical coefficients represent one-standard-deviation changes.

\begin{lstlisting}[style=mlpython]
linear = fitted["linear regression"]
names = linear.named_steps["preprocess"].get_feature_names_out()
coefs = pd.Series(linear.named_steps["model"].coef_, index=names).sort_values()

plt.figure(figsize=(6, 3.6))
plt.barh(coefs.index, coefs.values)
plt.axvline(0, color="grey", linewidth=0.8)
plt.xlabel("Change in predicted charges (dollars)")
plt.title("What the linear model learned")
plt.tight_layout()
plt.show()
\end{lstlisting}

\textbf{Inspect.} The smoker coefficient exceeds \$20,000; age and BMI contribute more than region. These are model associations, not causal effects.

Plot ridge residuals; structure suggests missed patterns.

\begin{lstlisting}[style=mlpython]
ridge_pred = predictions["ridge"]
residual = y_valid.to_numpy() - ridge_pred
is_smoker = (X_valid["smoker"] == "yes").to_numpy()

fig, axes = plt.subplots(1, 2, figsize=(9.5, 3.8))
axes[0].scatter(y_valid[~is_smoker], ridge_pred[~is_smoker], alpha=0.5, label="non-smoker")
axes[0].scatter(y_valid[is_smoker], ridge_pred[is_smoker], alpha=0.5, label="smoker")
axes[0].plot([0, y_valid.max()], [0, y_valid.max()], linestyle="--", color="grey")
axes[0].set_xlabel("Actual charges")
axes[0].set_ylabel("Predicted charges")
axes[0].set_title("Ridge: actual versus predicted")
axes[0].legend()
axes[1].scatter(ridge_pred, residual, alpha=0.6)
axes[1].axhline(0, linestyle="--", color="grey")
axes[1].set_xlabel("Predicted charges")
axes[1].set_ylabel("Residual (actual - predicted)")
axes[1].set_title("Ridge: residual plot")
plt.tight_layout()
plt.show()
\end{lstlisting}

\textbf{Inspect.} Smoker bands suggest a smoker $\times$ BMI interaction. Shrinkage alone cannot add this missing structure.

\medskip\noindent\textbf{Part B --- Regression on a time series: validate in date order.}\par

Sort two years of daily rentals by date; reserve the last 20\% for validation to avoid training on future days.

\begin{lstlisting}[style=mlpython]
bikes = pd.read_csv(DATA / "bike_sharing_dataset_day.csv")
bikes["dteday"] = pd.to_datetime(bikes["dteday"])
bikes = bikes.sort_values("dteday").reset_index(drop=True)

# cnt = casual + registered, so those two columns would leak the answer.
bike_cat = ["season", "yr", "mnth", "holiday", "weekday", "workingday", "weathersit"]
bike_num = ["temp", "atemp", "hum", "windspeed"]
X_b = bikes[bike_cat + bike_num]
y_b = bikes["cnt"]

cut = int(len(bikes) * 0.80)
print("train:", bikes["dteday"].iloc[0].date(), "to", bikes["dteday"].iloc[cut - 1].date())
print("valid:", bikes["dteday"].iloc[cut].date(), "to", bikes["dteday"].iloc[-1].date())
\end{lstlisting}

One-hot encode calendar categories; pass numerical weather features directly to the unscaled tree model.

\begin{lstlisting}[style=mlpython]
Xb_train, Xb_valid = X_b.iloc[:cut], X_b.iloc[cut:]
yb_train, yb_valid = y_b.iloc[:cut], y_b.iloc[cut:]
dates_valid = bikes["dteday"].iloc[cut:]

bike_model = Pipeline([
    ("preprocess", ColumnTransformer([
        ("cat", OneHotEncoder(handle_unknown="ignore"), bike_cat),
        ("num", "passthrough", bike_num),
    ])),
    ("model", RandomForestRegressor(
        n_estimators=400, min_samples_leaf=2, random_state=42, n_jobs=-1
    )),
])
bike_model.fit(Xb_train, yb_train)
bike_pred = bike_model.predict(Xb_valid)
print("daily MAE:", round(mean_absolute_error(yb_valid, bike_pred), 1),
      " (mean daily demand in validation:", round(yb_valid.mean(), 1), ")")

plt.figure(figsize=(8, 3.6))
plt.plot(dates_valid, yb_valid.to_numpy(), label="actual")
plt.plot(dates_valid, bike_pred, label="predicted")
plt.xlabel("Date")
plt.ylabel("Daily rentals")
plt.title("Chronological validation: actual versus predicted demand")
plt.legend()
plt.tight_layout()
plt.show()
\end{lstlisting}

\textbf{Inspect.} The forest tracks seasonality but underpredicts peaks beyond the training-label range; it cannot extrapolate.

Repeat chronologically on 17,379 hourly rows, adding \texttt{hr}; plot one validation week.

\begin{lstlisting}[style=mlpython]
hourly = pd.read_csv(DATA / "bike_sharing_dataset_hour.csv")
hourly["dteday"] = pd.to_datetime(hourly["dteday"])
hourly = hourly.sort_values(["dteday", "hr"]).reset_index(drop=True)

hour_cat = bike_cat + ["hr"]
Xh = hourly[hour_cat + bike_num]
yh = hourly["cnt"]
cut_h = int(len(hourly) * 0.80)

hour_model = Pipeline([
    ("preprocess", ColumnTransformer([
        ("cat", OneHotEncoder(handle_unknown="ignore"), hour_cat),
        ("num", "passthrough", bike_num),
    ])),
    ("model", RandomForestRegressor(
        n_estimators=200, min_samples_leaf=2, random_state=42, n_jobs=-1
    )),
])
hour_model.fit(Xh.iloc[:cut_h], yh.iloc[:cut_h])
hour_pred = hour_model.predict(Xh.iloc[cut_h:])
print("hourly MAE:", round(mean_absolute_error(yh.iloc[cut_h:], hour_pred), 1),
      " (mean hourly demand in validation:", round(yh.iloc[cut_h:].mean(), 1), ")")

week = slice(0, 24 * 7)
plt.figure(figsize=(9, 3.4))
plt.plot(range(24 * 7), yh.iloc[cut_h:].to_numpy()[week], label="actual")
plt.plot(range(24 * 7), hour_pred[week], label="predicted")
plt.xlabel("Hour of the first validation week")
plt.ylabel("Rentals per hour")
plt.title("Hourly demand has a shape that hr captures")
plt.legend()
plt.tight_layout()
plt.show()
\end{lstlisting}

\textbf{Inspect.} Working days have two peaks; weekends one. Lab~8 predicts this series from 24 previous hours.

\medskip\noindent\textbf{Part C --- Classification: baselines, probabilities, and the threshold.}\par

Compare logistic regression and a forest using seven passenger features, no identifiers, and stratified splits. Accuracy evaluates thresholded decisions; AUC evaluates ranking.

\begin{lstlisting}[style=mlpython]
titanic = pd.read_csv(DATA / "titanic_dataset.csv")

features = ["Pclass", "Sex", "Age", "SibSp", "Parch", "Fare", "Embarked"]
t_num = ["Age", "SibSp", "Parch", "Fare"]
t_cat = ["Pclass", "Sex", "Embarked"]
X_t = titanic[features]
y_t = titanic["Survived"]

Xt_train, Xt_valid, yt_train, yt_valid = train_test_split(
    X_t, y_t, test_size=0.25, random_state=42, stratify=y_t
)
print("majority-class baseline accuracy:", round(1 - yt_valid.mean(), 3))

scores = {}
fitted_t = {}
for name, estimator in {
    "logistic": LogisticRegression(max_iter=1000),
    "random_forest": RandomForestClassifier(n_estimators=300, min_samples_leaf=2,
                                            random_state=42),
}.items():
    pipe = Pipeline([("preprocess", make_preprocess(t_num, t_cat)), ("model", estimator)])
    pipe.fit(Xt_train, yt_train)
    fitted_t[name] = pipe
    scores[name] = {"accuracy": accuracy_score(yt_valid, pipe.predict(Xt_valid)),
                    "roc_auc": roc_auc_score(yt_valid, pipe.predict_proba(Xt_valid)[:, 1])}
    print(f"{name:14s} accuracy={scores[name]['accuracy']:.3f} "
          f"roc_auc={scores[name]['roc_auc']:.3f}")

# The same logistic model, three decision rules (the table in the notes).
prob_t = fitted_t["logistic"].predict_proba(Xt_valid)[:, 1]
for threshold in [0.30, 0.50, 0.70]:
    pred_t = (prob_t >= threshold).astype(int)
    print(f"threshold {threshold:.2f}: precision={precision_score(yt_valid, pred_t):.2f} "
          f"recall={recall_score(yt_valid, pred_t):.2f}")

pd.DataFrame(scores).T.plot.bar(figsize=(5, 3.2), ylim=(0.5, 1.0), rot=0)
plt.axhline(1 - yt_valid.mean(), linestyle="--", color="grey", label="majority baseline")
plt.title("Two models, two metrics, one validation split")
plt.legend(loc="lower right")
plt.tight_layout()
plt.show()
\end{lstlisting}

\textbf{Inspect.} Both beat the baseline, but metric rankings differ. Use Part~D cross-validation to assess small differences.

For 30 tumor measurements, predict \texttt{M} versus \texttt{B}. Plot probabilities by true class; a threshold marks malignant decisions.

\begin{lstlisting}[style=mlpython]
cancer = pd.read_csv(DATA / "breast_cancer_dataset.csv")

X_c = cancer.drop(columns=["id", "diagnosis", "Unnamed: 32"], errors="ignore")
y_c = (cancer["diagnosis"] == "M").astype(int)

Xc_train, Xc_valid, yc_train, yc_valid = train_test_split(
    X_c, y_c, test_size=0.25, random_state=42, stratify=y_c
)

cancer_model = Pipeline([
    ("scale", StandardScaler()),
    ("model", LogisticRegression(max_iter=2000)),
])
cancer_model.fit(Xc_train, yc_train)
prob = cancer_model.predict_proba(Xc_valid)[:, 1]

plt.figure(figsize=(6, 3.4))
plt.hist(prob[yc_valid == 0], bins=25, alpha=0.6, label="benign (actual)")
plt.hist(prob[yc_valid == 1], bins=25, alpha=0.6, label="malignant (actual)")
plt.axvline(0.50, color="black", linestyle="--", label="threshold 0.50")
plt.axvline(0.30, color="red", linestyle=":", label="threshold 0.30")
plt.xlabel("Predicted probability of malignant")
plt.ylabel("Tumours")
plt.title("A threshold is a vertical line through this picture")
plt.legend()
plt.tight_layout()
plt.show()
\end{lstlisting}

Confusion-matrix rows are truth; columns are decisions. Precision measures correct malignant alerts; recall measures detected malignancies. Lowering the threshold increases alerts and cannot reduce recall.

\begin{lstlisting}[style=mlpython]
for threshold in [0.50, 0.30]:
    pred = (prob >= threshold).astype(int)
    print(f"\nthreshold = {threshold:.2f}")
    print(confusion_matrix(yc_valid, pred))
    print(classification_report(yc_valid, pred, digits=3))
    ConfusionMatrixDisplay.from_predictions(yc_valid, pred, display_labels=["benign", "malignant"])
    plt.title(f"Confusion matrix at threshold {threshold:.2f}")
    plt.tight_layout()
    plt.show()

thresholds = np.linspace(0.05, 0.95, 19)
precisions = [precision_score(yc_valid, (prob >= t).astype(int), zero_division=0)
              for t in thresholds]
recalls = [recall_score(yc_valid, (prob >= t).astype(int)) for t in thresholds]

plt.figure(figsize=(6, 3.4))
plt.plot(thresholds, precisions, marker="o", label="precision")
plt.plot(thresholds, recalls, marker="o", label="recall")
plt.xlabel("Decision threshold")
plt.ylabel("Score")
plt.title("The threshold trades recall for precision")
plt.legend()
plt.tight_layout()
plt.show()
\end{lstlisting}

ROC AUC measures ranking across classes; average precision summarizes precision--recall, emphasizing positive retrieval.

\begin{lstlisting}[style=mlpython]
print("ROC AUC:", round(roc_auc_score(yc_valid, prob), 3))
print("PR AUC:", round(average_precision_score(yc_valid, prob), 3))

fig, axes = plt.subplots(1, 2, figsize=(9, 3.6))
RocCurveDisplay.from_predictions(yc_valid, prob, ax=axes[0])
axes[0].set_title("ROC curve")
PrecisionRecallDisplay.from_predictions(yc_valid, prob, ax=axes[1])
axes[1].set_title("Precision-recall curve")
plt.tight_layout()
plt.show()
\end{lstlisting}

\medskip\noindent\textbf{Part D --- Tune with cross-validation, then test once.}\par

Reserve Titanic test rows. Within training data, use five-fold CV to tune \texttt{C} and \texttt{class\_weight}; refit the winner on all training rows.

\begin{lstlisting}[style=mlpython]
X_train, X_test, y_train, y_test = train_test_split(
    X_t, y_t, test_size=0.25, random_state=42, stratify=y_t
)

search = GridSearchCV(
    Pipeline([("preprocess", make_preprocess(t_num, t_cat)),
              ("model", LogisticRegression(max_iter=2000))]),
    param_grid={
        "model__C": [0.01, 0.1, 1.0, 10.0, 100.0],
        "model__class_weight": [None, "balanced"],
    },
    scoring="roc_auc",
    cv=StratifiedKFold(n_splits=5, shuffle=True, random_state=42),
    n_jobs=-1,
)
search.fit(X_train, y_train)

cv_table = pd.DataFrame(search.cv_results_)
plt.figure(figsize=(5.5, 3.4))
for weight, group in cv_table.groupby(cv_table["param_model__class_weight"].astype(str)):
    plt.errorbar(group["param_model__C"].astype(float), group["mean_test_score"],
                 yerr=group["std_test_score"], marker="o", capsize=3, label=f"class_weight={weight}")
plt.xscale("log")
plt.xlabel("C (weaker regularization to the right)")
plt.ylabel("Cross-validated ROC AUC")
plt.title("What the grid search saw")
plt.legend()
plt.tight_layout()
plt.show()
\end{lstlisting}

\textbf{Inspect.} Fold spread exceeds many differences between settings; avoid overinterpreting small gains.

Only now is the test set touched, once.

\begin{lstlisting}[style=mlpython]
test_prob = search.predict_proba(X_test)[:, 1]
test_auc = roc_auc_score(y_test, test_prob)

print("best parameters:", search.best_params_)
print("cross-validated ROC AUC:", round(search.best_score_, 3))
print("final test ROC AUC:", round(test_auc, 3))
\end{lstlisting}

\medskip\noindent\textbf{Part E --- Save the models and use them.}\par

Save preprocessing and estimator together with \texttt{joblib.dump}; add a card documenting inputs and evaluation.

\begin{lstlisting}[style=mlpython]
final_model = search.best_estimator_
model_path = MODELS / "titanic_survival.joblib"
joblib.dump(final_model, model_path)

card = {
    "task": "predict Titanic survival (1 = survived)",
    "features": features,
    "model": "logistic regression pipeline",
    "best_parameters": search.best_params_,
    "training_rows": int(len(X_train)),
    "final_test_roc_auc": round(float(test_auc), 3),
    "decision_threshold": 0.5,
    "scikit_learn_version": sklearn.__version__,
}
(MODELS / "titanic_survival_card.json").write_text(json.dumps(card, indent=2))
size_kb = model_path.stat().st_size / 1024
print(f"saved: {model_path.name} ({size_kb:.1f} KB) and titanic_survival_card.json")
\end{lstlisting}

Load the pipeline and predict new rows with matching columns, without retraining. Repeat for insurance customers differing in one feature.

\begin{lstlisting}[style=mlpython]
loaded = joblib.load(model_path)
new_passengers = pd.DataFrame([
    {"Pclass": 1, "Sex": "female", "Age": 29, "SibSp": 0, "Parch": 0, "Fare": 80.0, "Embarked": "C"},
    {"Pclass": 3, "Sex": "male",   "Age": 40, "SibSp": 0, "Parch": 0, "Fare":  7.9, "Embarked": "S"},
])
new_passengers["survival_probability"] = loaded.predict_proba(new_passengers)[:, 1].round(3)
print(new_passengers)

joblib.dump(fitted["linear regression"], MODELS / "insurance_charges.joblib")
pricing = joblib.load(MODELS / "insurance_charges.joblib")
new_customer = pd.DataFrame([
    {"age": 35, "sex": "female", "bmi": 27.5, "children": 2, "smoker": "no", "region": "southeast"},
    {"age": 35, "sex": "female", "bmi": 27.5, "children": 2, "smoker": "yes", "region": "southeast"},
])
new_customer["predicted_charges"] = pricing.predict(new_customer).round(0)
print(new_customer)
\end{lstlisting}

Run the saved passenger model from the terminal:

\begin{lstlisting}[style=mlterminal]
python predict_passenger.py --pclass 2 --sex female --age 30 --fare 20
\end{lstlisting}

\begin{tryit}
Run the Part~C classification workflow on \texttt{heart\_failure\_dataset.csv} (label \texttt{HeartDisease}). Build the column lists from the data types with \texttt{select\_dtypes}, reuse \texttt{make\_preprocess}, and report accuracy, ROC AUC, and the confusion matrix at threshold 0.5.
\end{tryit}

\begin{decisionmemo}
Choose a threshold when missed positives cost more than false alarms. Explain confusion-matrix changes and justify the choice using Part C validation evidence.
\end{decisionmemo}
\end{labproject}
